\documentclass[11pt,a4paper]{article}
\usepackage[margin=24mm]{geometry}
\usepackage{amsmath,amssymb,amsthm,booktabs,hyperref}
\setlength{\emergencystretch}{2em}
\newtheorem{theorem}{Theorem}
\newtheorem{lemma}[theorem]{Lemma}
\newcommand{\ii}{i_2}
\newcommand{\Aut}{\operatorname{Aut}}
\title{Counting involutions that generate a chiral polyhedron}
\author{Human Compression project}\date{}
\begin{document}\maketitle
We give one counting argument for fourteen sporadic cases of the existing
existence theorem of Leemans and Liebeck~\cite{LL}. It uses ordinary class
sizes and a complete list of maximal subgroup classes. Both are retained
classification inputs, with their exact numerical use recorded. The
same involution bound handles all fourteen cases without distinguishing
the conjugacy classes of reversing automorphisms.

\section{The common counting theorem}
For a finite group $H$, let $i_d(H)$ be the number of elements of exact
order $d$. Let $G$ be nonabelian simple and identify $G$ with its inner
automorphism group in $A=\Aut(G)$. Fix $x\in G$ of odd prime order $p$.
Choose one representative of each $G$-conjugacy class of maximal
subgroups, and define
\[
 B_G(x)=|C_G(x)|\sum_{M}\frac{i_p(M)\ii(M)}{|M|}.             \tag{1}
\]
Let $\mu$ be an upper bound for $\ii(C_G(\alpha))$ for every
involution $\alpha\in A$ that inverts $x$. Put
\[
 R_A(x)=
 \begin{cases}
 0,&x\text{ is not conjugate to }x^{-1}\text{ in }A,\\
 |C_A(x)|\mu,&\text{otherwise}.
 \end{cases}                                                \tag{2}
\]
The upper bound in (1) may be rational; it need not count an actual set
exactly.

\begin{theorem}\label{count}
If $\ii(G)>B_G(x)+R_A(x)$, there is an involution $t\in G$ such that
$G=\langle x,t\rangle$ and no automorphism of $G$ inverts $x$ while
fixing $t$.
\end{theorem}
\begin{proof}
Every maximal subgroup $M<G$ is self-normalizing. Indeed its normalizer
is $M$ or $G$, and the latter would make $M$ a proper nontrivial normal
subgroup of the simple group $G$. Count the conjugates of $M$ containing
$x$. The set of conjugating elements taking $x$ into $M$ has size
$|C_G(x)|\,|x^G\cap M|$. Each conjugate of $M$ is counted $|M|$
times. Hence the number of these conjugates is
\[
 \frac{|C_G(x)|\,|x^G\cap M|}{|M|}
       \le \frac{|C_G(x)|i_p(M)}{|M|}.
\]
If $\langle x,t\rangle$ is proper, it is contained in one of these
maximal subgroups. Summing their involution counts proves that at most
$B_G(x)$ involutions fail to generate with $x$.

If an automorphism $\alpha$ inverts $x$ and fixes a generating partner
$t$, then $\alpha^2$ fixes both generators and is the identity. Since
$p$ is odd, $\alpha$ is nonidentity and thus an involution. All elements
of $A$ taking $x$ to $x^{-1}$, if any exist, form a coset of $C_A(x)$:
the quotient of any two such elements centralizes $x$. There are at
most $|C_A(x)|$ possible inverting involutions. Each fixes at most
$\mu$ involutions of $G$. The union of these fixed sets has size at
most $R_A(x)$; it is empty if $x$ is not conjugate to its inverse in $A$.
The strict inequality therefore leaves an involution outside both
excluded sets. It has the required properties.
\end{proof}

The maximal-subgroup count and inverter coset are standard counting
arguments, closely related to the preliminary lemmas of~\cite{LL}.
Our common specialization uses a deliberately coarse bound for $\mu$.
It avoids a separate analysis of the reversing involutions in each group.

\begin{lemma}\label{mu}
If $A=G$, one may take
\[
             \mu=\max_{M\text{ maximal in }G}\ii(M).          \tag{3}
\]
Suppose instead that $[A:G]=2$ and
$|C_A(a)|=2|C_G(a)|$ for every involution $a\in G$.
Then one may take
\[
        \mu=\max_{M\text{ maximal in }A,\ M\ne G}\ii(M).     \tag{4}
\]
\end{lemma}
\begin{proof}
The centre of $A$ is trivial. An automorphism central in $A$ commutes
with every inner automorphism; therefore it fixes every element of
the centreless group $G$ and is the identity. In particular the
centralizer of a nonidentity involution in $A$ is proper and lies in
a maximal subgroup. This proves (3), since
$C_G(\alpha)\le C_A(\alpha)$.

For (4), if $\alpha\notin G$, its centralizer contains $\alpha$
and is not contained in $G$. If $\alpha\in G$, the stated centralizer
index likewise shows that $C_A(\alpha)$ contains an element outside
$G$. A maximal subgroup containing this centralizer therefore cannot
be $G$. Containment bounds its involution count by (4), and again
$C_G(\alpha)\le C_A(\alpha)$.
\end{proof}

The centralizer condition in the second part is checked for every
involution class of $G$, using its actual fusion into $A$. It is not
inferred merely from the fact that $A/G$ has order two.

\section{Fourteen applications}
For each group below choose the specified prime-order class in the
bound data. The column $L$ is a lower bound for $\ii(G)$, and $B,R$
are the exact upper bounds computed from (1)--(4). In every row
$B+R<L$. Only for $B$ and $M$ do we replace the exact involution count
by a smaller power of ten in this display; the full counts remain in
the certificate.
\begin{center}\small
\begin{tabular}{lrrrr}\toprule
$G$&$p$&$L$&$B$&$R$\\\midrule
$J_4$&29&51747149311&29&101486051\\
$\mathrm{Fi}_{23}$&23&12895196031&976120&0\\
$\mathrm{Fi}'_{24}$&29&7824165773823&58&747921369827\\
$M$&71&$10^{20}$&4970&0\\
$B$&47&$10^{12}$&0&0\\
$Th$&31&976841775&44826&0\\
$Ly$&37&1296826875&74&15054375\\
$Co_1$&23&12792894975&72100538&0\\
$Co_2$&23&29771775&7590&0\\
$Co_3$&23&2778975&7590&0\\
$Ru$&29&1846575&870&390195\\
$\mathrm{Fi}_{22}$&13&37706175&1443546&18669235\\
$\mathrm{O'N}$&31&2857239&1860&95945\\
$HN$&19&75603375&7182&2662869\\\bottomrule
\end{tabular}
\end{center}
The full automorphism groups equal $G$ except in the four rows
$\mathrm{Fi}'_{24},\mathrm{Fi}_{22},\mathrm{O'N},HN$, where they
are the retained extensions of index two. The ordinary tables and
their actual class fusions verify the condition in Lemma~\ref{mu}
in each of those four cases. For the rows with $R=0$, the inverse
class differs from the chosen class in the full automorphism group,
not merely in an unspecified subgroup.

For clarity, the nonzero values of $\mu$ actually used are
\begin{center}
\begin{tabular}{lrlr}\toprule
$G$&$\mu$&$G$&$\mu$\\\midrule
$J_4$&3499519&$\mathrm{Fi}_{22}$&1436095\\
$Ly$&406875&$\mathrm{O'N}$&3095\\
$\mathrm{Fi}'_{24}$&25790392063&$HN$&140151\\
$Ru$&13455&&\\\bottomrule
\end{tabular}
\end{center}
These are maxima over complete maximal-subgroup lists, as specified
in (3) or (4). For example, the largest value for the automorphism
group of $\mathrm{Fi}'_{24}$, after omitting its simple socle, is
$i_2(2\times\mathrm{Fi}_{23})=25790392063$. Thus even this coarse
bound suffices. No assertion that an inverter belongs to a particular
class of involutions is needed.

Theorem~\ref{count} gives the required generating pairs. For the
polyhedron conclusion put $y=x^{-1}t$. The cyclic intersection
$\langle x\rangle\cap\langle y\rangle$ is central in
$\langle x,y\rangle=G$, so it is trivial. Also $(xy)^2=1$.
The element $y$ cannot be an involution: otherwise $G$ would be
generated by the two involutions $y,t$ and would be dihedral.
The case $y=1$ would make $G$ cyclic. Hence $|x|,|y|>2$.
The retained Schulte--Weiss reconstruction theorem supplies a chiral
polyhedron with full automorphism group $G$, since the reversing
automorphism has been excluded.

\section{Inputs, verification, and attribution}
The source's existing existence conclusions are being reproved, not
claimed as new. The retained facts are the identifications of the
sporadic groups and their automorphism groups, the complete lists
of maximal subgroup classes, the ordinary class sizes and element
orders, and the displayed inclusion fusions. The verifier obtains
the numerical data from GAP's Character Table Library, in its
recorded edition. Its \texttt{Maxes} attribute is documented to give
one representative for every conjugacy class of maximal subgroups
when it is present. Completeness of those classification lists is
a premise; it is not proved by iterating over a software list.

For every required table, the checker sums the sizes of the classes
of exact order two or $p$. It retains every maximal-subgroup class,
including repetitions of a table name when distinct subgroup classes
have the same character table. It verifies each centralizer index
used in (4), the actual inverse-class test in (2), and all fourteen
strict inequalities with exact integers or rationals. The full
underlying rows and their source bindings are retained, so the small
displayed table is not an unexplained set of trusted totals.

Together with the separately reviewed twelve explicit word
certificates, these fourteen cases cover all twenty-six sporadic
groups in the source theorem. The alternating branch has its own
uniform proof. The classical and exceptional Lie families remain
separate obligations. 
\begin{thebibliography}{9}
\bibitem{LL} D.~Leemans and M.~W.~Liebeck, \emph{Chiral polyhedra and
finite simple groups}, fixed arXiv:1603.07713v3, Sections~2 and~6.
\bibitem{SW} E.~Schulte and A.~Ivi\'c Weiss, \emph{Chiral polytopes},
Applied Geometry and Discrete Mathematics, DIMACS Series 4 (1991),493--516.
\bibitem{Atlas} J.~H.~Conway, R.~T.~Curtis, S.~P.~Norton, R.~A.~Parker
and R.~A.~Wilson, \emph{Atlas of Finite Groups}, Oxford University Press, 1985.
\bibitem{CTblLib} T.~Breuer, \emph{The GAP Character Table Library},
version 1.3.11, data and manual Section~3.7.1; the exact retained
installation files are specified in the dependency ledger.
\end{thebibliography}
\end{document}
